% Idealized FLFM light path (supervisor, 2026-11-23 meeting): the microlens
% array sits directly in the back pupil of the objective and the sensor in the
% lenslets' focal plane, with no relays. The full path with both relays is
% flfm-full-path.tex, shown at the end of the chapter. Five lenslets drawn
% (N = 5); not to scale. Replaces flfm-path.tex as Figure 3.1; that file is
% kept only for Appendix D.
\begin{tikzpicture}[
    x=1cm,y=1cm,
    font=\footnotesize,
    ink/.style={text=black!85},
    lens/.style={<->,>={Stealth[length=2mm]},draw=black!75,line width=0.8pt},
    ray/.style={draw=blue!60!black,line width=0.7pt},
    slice/.style={draw=black!25,line width=0.4pt},
    plane/.style={draw=black!45,dashed,line width=0.5pt},
    subimg/.style={draw=orange!75!black,line width=2.2pt,opacity=0.55},
    dim/.style={<->,>={Stealth[length=1.6mm]},draw=black!70,line width=0.5pt},
    lead/.style={draw=black!55,line width=0.4pt},
    lab/.style={font=\footnotesize,ink,inner sep=1.5pt,align=center}
]
% ---- geometry (cm; not to scale) ----
% object 0 | objective 1.8 | back pupil = MLA 3.6 | sensor = MLA focal plane 5.6
\def\xo{1.8}  \def\xp{3.6}  \def\xs{5.6}
\def\dm{0.56}                           % lenslet pitch d_MLA
\def\hp{1.4}                            % half of the pupil (5 pitches)
\def\sfr{0.8}                           % image of the field, fraction of one lenslet pitch

% optical axis
\draw[draw=black!30,line width=0.4pt] (-0.4,0) -- (\xs+0.5,0);

% object plane
\draw[plane] (0,-1.6) -- (0,1.6);

% objective
\draw[lens] (\xo,-1.6) -- (\xo,1.6);

% microlens array in the back pupil: 5 lenslets
\foreach \k in {-2,...,2} {
    \draw[draw=black!70,fill=black!6,line width=0.5pt]
        (\xp-0.07,{\k*\dm-0.5*\dm}) rectangle (\xp+0.07,{\k*\dm+0.5*\dm});
}

% rays of one on-axis emitter: the objective collimates, each lenslet focuses
\draw[ray] (0,0) -- (\xo,\hp);
\draw[ray] (0,0) -- (\xo,-\hp);
\draw[ray] (\xo,\hp) -- (\xp,\hp);
\draw[ray] (\xo,-\hp) -- (\xp,-\hp);
\foreach \k in {-2,...,2} {
    \draw[slice] (\xo,{\k*\dm-0.5*\dm}) -- (\xp,{\k*\dm-0.5*\dm});
    \draw[ray] (\xp,{\k*\dm-0.5*\dm}) -- (\xs,{\k*\dm});
    \draw[ray] (\xp,{\k*\dm+0.5*\dm}) -- (\xs,{\k*\dm});
}
\draw[slice] (\xo,{2.5*\dm}) -- (\xp,{2.5*\dm});

% sensor in the lenslets' focal plane, with pixel ticks
\draw[draw=black!70,fill=black!12,line width=0.5pt] (\xs,-1.45) rectangle (\xs+0.12,1.45);
\foreach \y in {-1.4,-1.345,...,1.42} {
    \draw[draw=black!45,line width=0.25pt] (\xs,\y) -- (\xs+0.12,\y);
}
% one view per lenslet: the image of the field (bar) and the emitter's spot
\foreach \k in {-2,...,2} {
    \draw[subimg] (\xs-0.03,{\k*\dm-0.5*\sfr*\dm}) -- (\xs-0.03,{\k*\dm+0.5*\sfr*\dm});
    \fill[orange!80!black] (\xs-0.03,{\k*\dm}) circle (0.045);
}

% dimensions at the pupil
\draw[dim] (\xp-0.75,-\hp) -- (\xp-0.75,\hp);
\node[lab,anchor=east,fill=white,inner sep=1pt] at (\xp-0.8,-0.55) {$D_{pupil}$};
\draw[dim] (\xp-0.35,{1.5*\dm}) -- (\xp-0.35,{2.5*\dm});
\node[lab,anchor=east,fill=white,inner sep=1pt] at (\xp-0.4,{2.0*\dm}) {$d_{MLA}$};

% focal lengths along the bottom
\def\yb{-2.3}
\draw[dim] (0,\yb) -- (\xo,\yb);    \node[lab,fill=white] at ({0.5*\xo},\yb) {$f_{obj}$};
\draw[dim] (\xo,\yb) -- (\xp,\yb);  \node[lab,fill=white] at ({0.5*(\xo+\xp)},\yb) {$f_{obj}$};
\draw[dim] (\xp,\yb) -- (\xs,\yb);  \node[lab,fill=white] at ({0.5*(\xp+\xs)},\yb) {$f_{MLA}$};

% names
\node[lab,anchor=south] at (0,1.65) {object\\plane};
\node[lab,anchor=south] at (\xo,1.65) {objective};
\node[lab,anchor=south] at (\xp,1.65) {MLA in the\\back pupil};
\node[lab,anchor=south] at (\xs+0.06,1.65) {sensor\\(MLA focal\\plane)};
\node[lab,anchor=north] at (\xp,-1.65) {$N\approx D_{pupil}/d_{MLA}$};

% ---- inset: one view on the sensor (zoom of the central view) ----
\def\xi{7.0}   \def\px{0.2}
\pgfmathsetmacro{\xiE}{\xi+10*\px}
\draw[draw=black!70,line width=0.6pt] (\xs+0.25,{-0.5*\dm}) -- (\xs+0.25,{0.5*\dm});
\draw[lead,dashed] (\xs+0.25,{0.5*\dm}) -- (\xi,1.0);
\draw[lead,dashed] (\xs+0.25,{-0.5*\dm}) -- (\xi,-1.0);
% image of the field within one lenslet pitch (no overlap with the neighbors)
\fill[orange!25] (\xi,{-1.0+\px}) rectangle (\xiE,{1.0-\px});
\draw[draw=black!25,line width=0.3pt] (\xi,-1.0) grid[step=\px] (\xiE,1.0);
\draw[draw=black!70,line width=0.5pt] (\xi,-1.0) rectangle (\xiE,1.0);
% one pixel, and the emitter's spot spanning alpha = 2 pixels
\fill[black!45] ({\xi+7*\px},{-1.0+2*\px}) rectangle ({\xi+8*\px},{-1.0+3*\px});
\fill[orange!80!black,opacity=0.6] ({\xi+5*\px},0) circle ({\px});
\draw[dim] (\xiE+0.15,{-1.0+\px}) -- (\xiE+0.15,{1.0-\px});
\node[lab,anchor=west] at (\xiE+0.2,0.25) {$FOV$};
\node[lab,anchor=south] at ({0.5*(\xi+\xiE)},1.05) {one view:\\$\approx N_{px}/N$ pixels};
\node[lab,anchor=north,align=left] at ({0.5*(\xi+\xiE)+0.3},-1.1)
    {gray: one pixel ($p$ in\\object space)\\spot: $\Delta x_{LF}=\alpha\,p$};
\end{tikzpicture}
